%subfiles使用時はdocument環境内にサブファイルも記入する
%サブファイル内のプリアンブルはmain.texのビルド時は無視される
\documentclass[../../main.tex]{subfiles} %[]内にこのドキュメントからmain.texまでのpathを記入

\graphicspath{{../../fig/}} %このドキュメントから画像ディレクトリまでのpathを記述
\setcounter{chapter}{0} %チャプター番号の調整

%***************************************** ＜Main body＞ ********************************************%
\begin{document}

\setcounter{page}{0} %ページ番号をゼロに設定
\pagenumbering{arabic} %本文用のページ番号付け

\chapter{モジュール}\label{chap:introduction}


\section{Module of Human–Robot Interaction with Correction and Suggestion (HRI-CAS)}
\subsection{Target Bottle Inference (TBI) Module}

\textbf{Definition of Features}:  
The TBI module aims to estimate the user’s intended target object based on bodily movements.  
Figure~3 illustrates the distribution of operational features in the information space.  
% Insert figure

Let the position of a target object $o_i$ be denoted as $\mathbf{p}_i(t)$,  
and the positions of the user’s eyes, head, 
and hand be $\mathbf{p}_{eye}(t)$, $\mathbf{p}_{head}(t)$, and $\mathbf{p}_{hand}(t)$, respectively.  
For each object, the normalized direction vectors from these body parts are defined as:
\begin{equation}
\mathbf{d}_i^{eye}(t) = \frac{\mathbf{p}_i(t)-\mathbf{p}_{eye}(t)}{\|\mathbf{p}_i(t)-\mathbf{p}_{eye}(t)\|}, \quad
\mathbf{d}_i^{head}(t) = \frac{\mathbf{p}_i(t)-\mathbf{p}_{head}(t)}{\|\mathbf{p}_i(t)-\mathbf{p}_{head}(t)\|}, \quad
\mathbf{d}_i^{hand}(t) = \frac{\mathbf{p}_i(t)-\mathbf{p}_{hand}(t)}{\|\mathbf{p}_i(t)-\mathbf{p}_{hand}(t)\|}.
\end{equation}
These vectors geometrically define the directions in which the user is focusing attention or exerting effort.  
To incorporate dynamics of motion, normalized gaze, head orientation, and hand acceleration vectors are introduced as  
$\hat{\mathbf{v}}_{gaze}(t)$, $\hat{\mathbf{v}}_{head}(t)$, and $\hat{\mathbf{a}}_{hand}(t)$.

\medskip
\textbf{Cosine Similarity Scores}:  
To evaluate whether the user’s body movements align with the direction of candidate objects,  
cosine similarity between these features and the direction vectors is computed:
\begin{align}
s_{\text{gaze}}(i,t) &= \max \left(0, \hat{\mathbf{v}}_{gaze}(t)\cdot \mathbf{d}_i^{eye}(t)\right), \\
s_{\text{head}}(i,t) &= \max \left(0, \hat{\mathbf{v}}_{head}(t)\cdot \mathbf{d}_i^{head}(t)\right), \\
s_{\text{accel}}(i,t) &= \max \left(0, \hat{\mathbf{a}}_{hand}(t)\cdot \mathbf{d}_i^{hand}(t)\right).
\end{align}
Here, $s_{\text{gaze}}$ represents the degree to which the gaze is directed toward each object,  
$s_{\text{head}}$ represents the degree to which the head orientation captures each object,  
and $s_{\text{accel}}$ represents the degree to which the hand acceleration is aligned with the object’s direction.  

In addition, physical graspability is considered by introducing:  
(i) a proximity score $s_{\text{prox}}$ based on the Euclidean distance between the hand and the object, and  
(ii) a binary contact variable $c_i(t)$:
\begin{equation}
s_{\text{prox}}(i,t) = \operatorname{clip}\left(1 - \frac{\|\mathbf{p}_{hand}(t)-\mathbf{p}_i(t)\|}{R_{\max}},\,0,\,1\right), \quad
c_i(t) \in \{0,1\}.
\end{equation}
Thus, the module integrates both body movement and physical graspability.

\medskip
\textbf{Score Integration}:  
The feature vector is defined as
\begin{equation}
\mathbf{s}(i,t) = [s_{\text{prox}},\,c_i,\,s_{\text{gaze}},\,s_{\text{head}},\,s_{\text{accel}}]^\top,
\end{equation}
and the integrated score is calculated with different weightings depending on whether 
the target belongs to the reachable set $\mathcal{B}_{in}(t)$  
or the unreachable set $\mathcal{B}_{out}(t)$:
\begin{equation}
S(i,t) =
\begin{cases}
\mathbf{w}_{in}^\top \mathbf{s}(i,t), & i \in \mathcal{B}_{in}(t), \\[6pt]
\mathbf{w}_{out}^\top \mathbf{s}(i,t), & i \in \mathcal{B}_{out}(t).
\end{cases}
\end{equation}
Here, $\mathbf{w}_{in}$ and $\mathbf{w}_{out}$ are predefined weight vectors.  
For reachable objects, emphasis is placed on contact and proximity, 
ensuring robustness against noisy gaze fluctuations.  
For unreachable objects, emphasis is placed on gaze and hand motion, capturing the user’s distal intention.  
This design enables flexible and error-tolerant inference depending on the situation.

\medskip
\textbf{Temporal Smoothing}:  
To suppress rapid fluctuations caused by momentary body movements,  
an exponential moving average is applied:
\begin{equation}
\bar{S}(i,t) = \alpha S(i,t) + (1-\alpha)\bar{S}(i,t-1), \quad \alpha \in (0,1].
\end{equation}
Here, $\bar{S}(i,t)$ is the smoothed integrated score, and $\alpha$ is the smoothing coefficient.  
This prevents transient misdetections and allows stable inference of user intention.

\medskip
\textbf{Final Decision Rule}:  
The final target object is selected as the one maximizing the smoothed score over the entire candidate set $\mathcal{B}$:
\begin{equation}
i^*(t) = \arg\max_{i \in \mathcal{B}} \bar{S}(i,t).
\end{equation}
However, the selected object must satisfy threshold conditions:
\begin{equation}
  \bar{S}\!\left(i^{*}(t),\,t\right) \;\geq\;
  \begin{cases}
    T_{\text{in}},  & i^{*}(t) \in \mathcal{B}_{\text{in}}(t), \\[6pt]
    T_{\text{out}}, & i^{*}(t) \in \mathcal{B}_{\text{out}}(t).
  \end{cases}
\end{equation}
Here, $T_{\text{in}}$ and $T_{\text{out}}$ are thresholds for the reachable and unreachable sets, respectively.  
This two-step decision ensures that only targets with sufficiently strong intentional signals are confirmed.

\medskip
\textbf{Information Presentation Adjustment}:  
When presenting inference results to the user,  
the TBI module adjusts the amount of information to prevent cognitive 
overload while supporting intuitive understanding.  

First, as shown in Figure~4, each candidate bottle is classified by 
the \textit{Reachability Map (RM)} as “graspable” or “non-graspable.”  
Graspable objects are displayed with a base color,  
while non-graspable ones are rendered semi-transparent with reduced saturation,  
ensuring they do not unnecessarily attract the user’s attention.  
% Insert figure

Furthermore, as illustrated in Figure~5,  
once the target is inferred, all non-selected objects are displayed in semi-transparent form.  
At this stage, different highlighting schemes are used for graspable and non-graspable objects:  
graspable targets are emphasized with a color indicating manipulability,  
while non-graspable targets are highlighted with a different tone to indicate the need for base relocation.  
% Insert figure

Through this two-step information adjustment,  
users can clearly distinguish between “objects that can currently be grasped” 
and “objects that require relocation.”  
By emphasizing the final inferred target,  
the system reduces selection errors and alleviates cognitive load,  
supporting intuitive decision-making.

\paragraph{Summary of TBI}  
In summary, the TBI module geometrically defines multiple bodily features 
such as gaze, head orientation, and hand motion,  
and integrates them through cosine similarity, proximity, and contact indicators 
to construct a robust scoring framework.  
Temporal smoothing suppresses transient misdetections, ensuring stable real-time inference.  

Furthermore, the inference results are visually presented in MR space:  
\begin{enumerate}
  \item Objects are classified by the Reachability Map (RM) as graspable or non-graspable,  
        displayed either in base color or semi-transparent with reduced saturation.  
  \item The inferred target bottle is emphasized,  
        with different highlight colors used for graspable and non-graspable objects,  
        providing explicit feedback to the user.  
\end{enumerate}

By jointly addressing “target inference” and “information presentation,”  
the TBI module reduces cognitive load, prevents misselection,  
and establishes accurate recognition of human intention as the foundation of the HRI-CAS framework.




\subsection{Base Relocation based on Inverse Reachability Map (BRI)}

When the target lies outside the robot’s reachable workspace, the robot base 
itself must be relocated.  
Naively exploring base candidates leads to prohibitive computational cost and 
makes it difficult 
to obtain configurations that respect human–robot collaboration.  
Therefore, as illustrated in Fig.~6, 
we introduce two concepts widely used in manipulator research— the \textit{Reachability Map (RM)} and
the \textit{Inverse Reachability Map (IRM)}—to systematically generate and evaluate candidate base poses \cite{ref6}. 
% insert paired figures

\medskip
\textbf{Basics of reachability and inverse reachability}:  
Given a robot base pose $\mathbf{b}$, the workspace that the end-effector (EE) can reach from $\mathbf{b}$ is defined as
\begin{equation}
RM(\mathbf{b}) = \{ \mathbf{p} \in \mathbb{R}^3 \mid \exists\, \theta \ \text{s.t.} \ FK(\theta, \mathbf{b}) = \mathbf{p} \},
\end{equation}
where $FK(\theta,\mathbf{b})$ denotes forward kinematics.  
In other words, if there exists a joint configuration $\theta$, the point $\mathbf{p}$ is reachable from the fixed base $\mathbf{b}$.  
Thus, the RM specifies “how far the arm can reach from a fixed base pose.”

Conversely, for a given task point $\mathbf{p}_t$, the set of base poses that makes $\mathbf{p}_t$ reachable is
\begin{equation}
\mathcal{B}(\mathbf{p}_t) = \{ \mathbf{b} \in \mathbb{R}^3 \mid \mathbf{p}_t \in RM(\mathbf{b}) \},
\end{equation}
which we call the \textit{Inverse Reachability Map (IRM)}.  
IRM is the inverse mapping that provides “the set of base poses from which the task point is reachable,” 
and is essential when base relocation is considered \cite{ref6}.

Beyond a binary reachable/unreachable decision, RM can evaluate stability/redundancy 
by counting the number 
of inverse kinematic (IK) solutions available at a point.  
In this work, the RM score for position $\mathbf{p}$ is defined by
\begin{equation}
S_{RM}(\mathbf{p}) = N_{\text{IK}}(\mathbf{p}).
\end{equation}
A larger number of solutions implies greater redundancy and flexibility in motion, 
which typically increases grasp success rates.

Since IRM is the inverse of RM, each candidate base pose $\mathbf{b}\in\mathcal{B}(\mathbf{p}_t)$ 
for target $\mathbf{p}_t$ inherits the corresponding RM score:
\begin{equation}
S_{IRM}(\mathbf{b}; \mathbf{p}_t) = N_{\text{IK}}(\mathbf{p}_t \mid \mathbf{b}).
\end{equation}
This allows us to assess candidate base configurations not only by “reachability,” 
but also by “how redundantly and stably” the target can be reached.

\medskip
\textbf{Optimization problem}:  
We now select the optimal base pose from the candidate set $\mathcal{B}(\mathbf{p}_t)$.  
As shown conceptually in Fig.~7, we formulate two cases depending on the availability of 
environment/navigation information.

\medskip
\textbf{Case: Known path}  
If a navigable path to the goal exists, we select the base by jointly considering 
the IRM score $S_{IRM}(\mathbf{b}; \mathbf{p}_t)$ and the movement cost $C(\mathbf{b})$:
\begin{equation}
\mathbf{b}^* = \arg\max_{\mathbf{b} \in \mathcal{B}(\mathbf{p}_t)}
\Big[\, S_{IRM}(\mathbf{b}; \mathbf{p}_t) - \lambda\, C(\mathbf{b}) \,\Big].
\end{equation}
Here, $C(\mathbf{b})$ denotes the cost to move from the current base $\mathbf{b}_0$ 
to $\mathbf{b}$ (e.g., Euclidean distance or shortest path length on an occupancy grid).  
The parameter $\lambda$ balances “reach stability” versus “movement efficiency.”  
The resulting candidate base poses and waypoint sequences are visualized in MR space; 
the user can then modify the path by directly manipulating the visualized waypoints.  
This mechanism realizes the coexistence of autonomous planning and user interaction, as depicted in Fig.~8.
% insert paired figures

\medskip
\textbf{Case: Unknown path}  
When environment information is incomplete and path generation fails, we adopt the IRT-MRO \cite{ref5} scheme:  
the user directly lays out a base route in MR space.  
On top of this, we introduce a local autonomous refinement using a cost map.  
As shown in Fig.~9, each waypoint is adjusted to a safe and feasible pose that accounts 
for obstacles and environmental changes.  
The refined route is immediately reflected in MR space so that the user can visually 
grasp the autonomous corrections made by the robot.
% (Mathematical details to be added)
% insert paired figures

\paragraph{Summary of BRI}  
In summary, our base relocation strategy adopts a two-track structure:  
(i) when a path is known, we perform autonomous optimization based on the IRM score and movement cost;  
(ii) when a path is unknown, we combine user-specified routes with cost-map-based refinement.  
This design satisfies both robot autonomy and user-driven interaction, achieving flexible and robust base relocation.



\subsection{Human-Intent Adaptive Grasp Manipulation (HIGM)}

Even when the base position is appropriate, discrepancies between human intent and the robot’s 
kinematic constraints can cause the end-effector (EE) to approach the target in unstable poses.  
This can lead to grasp failure or excessive joint load.  
To address this issue, we propose an \textit{adaptive grasp framework} that integrates Pregrasp design, 
null-space optimization, and a finite state machine (FSM).

\medskip
\textbf{Pregrasp and EE interpolation}:  
Human operation alone may command poses that are difficult for the robot to execute.  
For example, if the user attempts to pinch the top of a bottle, 
directly mapping the hand-back pose to the EE could result in an unstable grasp configuration.

To address this, as shown in Fig.~10, we generate a robot-friendly 
reference pose (Pregrasp pose $\mathbf{x}_{pre}$) based on the target position $\mathbf{p}_t$.  
The actual EE target is defined by interpolating between the user’s palm pose $EE_{palm}$ and 
the Pregrasp pose $EE_{pre}$, depending on the distance $d$ to the target:
\begin{equation}
EE_{target}(d) = (1-\beta(d))\,EE_{palm} + \beta(d)\,EE_{pre},
\end{equation}
where $\beta(d)\in[0,1]$ is a distance-dependent interpolation coefficient.

With this scheme, user intent ($EE_{palm}$) dominates when the target is far away, and as the robot approaches, 
the weight gradually shifts toward the stable Pregrasp pose ($EE_{pre}$).  
In other words, the user’s motion is “magnetically guided” toward the 
object—similar to aim assist in video games—preserving 
intuitive control while ensuring stable and safe grasping for the robot.

Furthermore, Pregrasp interpolation provides the global reference, 
while null-space optimization ensures local stability and safety.  
Thus, interpolation defines “which pose should be aimed for,” 
and null-space optimization ensures “how to safely approach that pose.”  
Combining them achieves grasping that respects user intent while adapting to the robot’s kinematic limits.  
% insert figure

\medskip
\textbf{Null-space optimization}:  
For redundant manipulators, multiple joint configurations can realize the same EE pose.  
The set of joint motions that do not affect the EE is called the \textit{null space} \cite{}.  

Formally, the task-space velocity $\dot{x}$ and joint velocity $\dot{\theta}$ satisfy
\begin{equation}
\dot{x} = J \dot{\theta}.
\end{equation}
A velocity $\dot{\theta}_{null}$ belongs to the null space if
\begin{equation}
J \dot{\theta}_{null} = 0.
\end{equation}
Such motions leave the EE unaffected while adjusting only the joint posture.

We exploit this property to simultaneously achieve the primary task and optimize joint posture.  
The joint velocity is given by
\begin{equation}
\dot{\theta} = J^+ \dot{x}_d + P \nabla_\theta H(\theta),
\end{equation}
where the damped least-squares pseudoinverse is
\begin{equation}
J^+ = J^\top (JJ^\top + \lambda I)^{-1},
\end{equation}
and $P = I - J^+ J$ is the null-space projector.  
The first term guarantees EE tracking, faithfully reproducing the user’s hand motion.  
The second term optimizes posture in the null space using the merit function $H(\theta)$ (Fig.~10):

\begin{equation}
H(\theta) = w_{mani}\log\det(JJ^\top) + w_{limit}\sum_i \log m_i(\theta) - \gamma\|\theta-\theta_{pre}\|^2.
\end{equation}

The roles of each term are:
\begin{itemize}
  \item $w_{mani}\log\det(JJ^\top)$: maximizes Yoshikawa’s manipulability, 
  avoiding singularities and preventing the arm from losing freedom.  
  This maintains motion redundancy, providing the user with smooth and continuous operation.

  \item $w_{limit}\sum_i \log m_i(\theta)$: enlarges the margin $m_i(\theta)$ to 
  joint limits, preventing joints from approaching boundaries.  
  This ensures mechanical safety and reduces the risk of sudden motion stops, improving user comfort.

  \item $-\gamma\|\theta-\theta_{pre}\|^2$: maintains consistency with the Pregrasp pose, 
  preventing optimization from drifting too far from user intent.  
  This balances intuitive operation and kinematic stability.
\end{itemize}

Thus, null-space optimization ensures  
(1) stability through singularity avoidance,  
(2) safety through joint-limit avoidance, and  
(3) intent preservation through Pregrasp consistency.  
It functions not only as numerical stabilization, but as a foundation linking 
“robotic physical constraints” with “the user’s continuous intuitive experience.”  
% insert figure

\medskip
\textbf{Finite State Machine (FSM)}:  
To manage grasp execution safely and reliably, we introduce a four-state FSM \cite{}.  
The FSM ensures both user-intent reflection and robot safety, preventing failures such 
as entering the target in unstable postures.

The roles and transitions of the states are:

\begin{enumerate}
  \item \textbf{TRACK}:  
  The robot follows the target using Palm/Pregrasp interpolation.  
  Mathematically, the EE target pose $EE_{target}(d)$ is generated 
  and joint angles are updated via DLS-IK.  
  User experience: “the robot hand smoothly follows my hand.”

  \item \textbf{BACKOFF}:  
  If manipulability $M(\theta)<\tau_{mani}$ or margin $m_i(\theta)<\tau_{limit}$, 
  the robot retreats to avoid instability.  
  This prevents “robot jamming” and provides the user with a sense of safety.

  \item \textbf{REALIGN}:  
  After BACKOFF, null-space optimization realigns the joints 
  until $M(\theta)\geq \tau_{mani}$ and $m_i(\theta)\geq \tau_{limit}$.  
  User experience: “the robot autonomously reconfigures itself,” revealing its adaptiveness.

  \item \textbf{APPROACH}:  
  Once stability is secured and the target distance $d<\tau_{dist}$, 
  the EE approaches to grasp.  
  This corresponds to “the robot decisively grabs,” improving grasp success.
\end{enumerate}

The FSM links mathematical criteria (manipulability, joint margin, distance thresholds) with 
user experience (tracking feel, safety, adaptiveness, visible grasp action).  
In particular, the BACKOFF $\to$ REALIGN transition guarantees that even if the user makes errors, 
the robot autonomously corrects them, realizing “collaborative safety” in HRC.

% insert figure

\subsection{Significance of Module Integration}

The proposed HRI-CAS consists of three modules: Target Bottle Inference (TBI), 
Base Relocation based on IRM (BRI), 
and Human-Intent Adaptive Grasp Manipulation (HIGM).  
Each module functions independently, but more importantly, 
they complement each other to overcome 
the limitations faced by conventional MR-based HRC systems.

First, TBI accurately estimates the user’s intent, clarifying “which object should be manipulated.”  
This object information is then passed to BRI, which determines whether base relocation is necessary 
depending on the object’s reachability.  
Subsequently, HIGM leverages the outputs of TBI and BRI to compensate for discrepancies between human 
intent and robot kinematic constraints, ensuring stable grasping.

Through the integration of TBI, BRI, and HIGM, the system establishes a coherent workflow:  
(1) recognizing user intent with high precision,  
(2) ensuring reachability through appropriate base relocation, and  
(3) guaranteeing grasp stability.  
As a result, the robot operates intuitively and adaptively without over-relying on the user’s physicality, 
providing a stable interaction experience for both novice and expert users.

Thus, the significance of this research lies not merely in individual technical improvements, but in the establishment 
of an integrated framework for HRC in MR environments, encompassing the flow of 
“intent recognition → reachability assurance → stable grasp.”


\end{document}